\documentclass[10pt]{beamer}

\usetheme{metropolis}
\definecolor{WiLabRed}{RGB}{197,18,48}
\setbeamercolor{frametitle}{fg=white,bg=WiLabRed}
\setbeamercolor{progress bar}{fg=WiLabRed!90}
\setbeamercolor{title separator}{fg=WiLabRed!90}
\setbeamercolor{progress bar in section page}{fg=WiLabRed!90}
\setbeamercolor{background canvas}{bg=white}
\setbeamercolor{alerted text}{fg=WiLabRed!90}

\usepackage{appendixnumberbeamer}

\usepackage{booktabs}
\usepackage[scale=2]{ccicons}

\usepackage{pgfplots}
\usepgfplotslibrary{dateplot}

\usepackage{xspace}
\newcommand{\themename}{\textbf{\textsc{metropolis}}\xspace}

\usepackage{marvosym}
%\usepackage{subfig}
\usepackage{graphicx}\graphicspath{{images/}}
\usepackage{subcaption}
\usepackage[framed]{./support/mcode}
\usepackage{listings}
\usepackage{amsmath}

\DeclareMathOperator{\sinc}{sinc}

\title{Laboratory 1 (Part 1): Square-Root Raised Cosine Pulse Shaping, Timing Error}
\subtitle{\textit{Software Defined Radio for Engineers} (Collins~\textit{et al.}), \textsection{6.1-6.2}}
\date{}
\author{\textbf{Alexander M. Wyglinski, Ph.D.}}
\institute{ \vspace*{1in}\hfill\includegraphics[height=1.125cm]{wilab_logo-A70916.eps} \qquad \includegraphics[height=1.125cm]{WPI_Inst_Prim_FulClr.eps}}
% \titlegraphic{\hfill\includegraphics[height=1.5cm]{logo.pdf}}

% Foot for all slides
\setbeamertemplate{frame footer}{\tiny \copyright~2018 by Alexander Wyglinski. This work is licensed under the Creative Commons Attribution-ShareAlike 4.0 International License. To view a copy of this license, visit http://creativecommons.org/licenses/by-sa/4.0/.}

\begin{document}

%\captionsetup[subfigure]{labelformat=empty}

%%%%%%%%%%%%%%%%%%%%%%%%%%%%%%%%%%%%%%%%%%%%%%%%%%%%%%%%%%

\maketitle




%%%%%%%%%%%%%%%%%%%%%%%%%%%%%%%%%%%%%%%%%%%%%%%%%%%%%%%%%%

\begin{frame}[fragile]{Square-Root Raised Cosine Pulse Shapes}


Square-root raised cosine (SRRC) has the impulse response:
%
\begin{equation}
	h_{\sqrt{\mathrm{RC}}}(t) = \begin{cases}
 \dfrac{1}{\sqrt{T_s}} \left( 1-\beta+4\dfrac{\beta}{\pi} \right),
       & t = 0 \\

\dfrac{\beta}{\sqrt{2T_s}}
\left[
\left(1+\dfrac{2}{\pi}\right)\sin\left(\dfrac{\pi}{4\beta}\right) +
\left(1-\dfrac{2}{\pi}\right)\cos\left(\dfrac{\pi}{4\beta}\right)
\right],
       & t = \pm \dfrac{T_s}{4\beta} \\

\dfrac{1}{\sqrt{T_s}} \dfrac{\sin\left[\pi \dfrac{t}{T_s}\left(1-\beta\right)\right] +
4\beta\dfrac{t}{T_s}\cos\left[\pi\dfrac{t}{T_s}\left(1+\beta\right)\right]}{\pi
\dfrac{t}{T_s}\left[1-\left(4\beta\dfrac{t}{T_s} \right)^2 \right]},
       & \mbox{otherwise}
\end{cases}
\end{equation}
where $T_s$ is the symbol period and $\beta \in \big[0,1\big]$ is the roll-off factor. 

\end{frame}


%%%%%%%%%%%%%%%%%%%%%%%%%%%%%%%%%%%%%%%%%%%%%%%%%%%%%%%%%%

\begin{frame}[fragile]{Why SRRCs?}

\centering
$H_{\sqrt{\mathrm{RC}}}(f)=\sqrt{H_{\mathrm{RC}}(f)}$\\
$\downarrow$\\
$H_{\mathrm{RC}}(f)$ satisfy Nyquist Criterion for Zero Intersymbol Interference (ISI)\\
$\downarrow$\\
Zero ISI results in transmissions whose symbols are not corrupted by previous and/or future symbols
 
\end{frame}


%%%%%%%%%%%%%%%%%%%%%%%%%%%%%%%%%%%%%%%%%%%%%%%%%%%%%%%%%%

\begin{frame}[fragile]{Raised Cosine Pulse Shapes}

We define a raised cosine (RC) pulse shape as:
\begin{equation}
 H_{\mathrm{RC}}(f)=
 \begin{cases}
  T,& 0\le{|f|}\le\frac{1-\beta}{2T}\\
  \frac{T}{2}\left(1+\cos{\left(\frac{\pi{t}}{\beta}\left(|f|-\frac{1-\beta}{T}\right)\right)}\right),&\frac{1-\beta}{2T}\le{|f|}\le\frac{1+\beta}{2T}\\
  0,&\mathrm{otherwise}
 \end{cases}
\end{equation}
\begin{equation}
 h_{\mathrm{RC}}=\frac{\cos(2\pi\beta{t}/T)\sinc(t/T)}{1-(2\beta{t}/T)^2}
\end{equation}

\end{frame}


%%%%%%%%%%%%%%%%%%%%%%%%%%%%%%%%%%%%%%%%%%%%%%%%%%%%%%%%%%

\begin{frame}[fragile]{Raised Cosine Pulse Shapes}

\centering

How do these pulse shapes look like in the time and frequency domains?
 
\end{frame}


%%%%%%%%%%%%%%%%%%%%%%%%%%%%%%%%%%%%%%%%%%%%%%%%%%%%%%%%%%

\begin{frame}[fragile]{Using SRRC Pulse Shapes}


Since SRRC pulse shapes are like any other type of filter, one can use a series of impulses defined by symbol data to produce a pulse shaped output

$\rightarrow$ \textit{Impulse Train Model}
 
\end{frame}


%%%%%%%%%%%%%%%%%%%%%%%%%%%%%%%%%%%%%%%%%%%%%%%%%%%%%%%%%%

\begin{frame}[fragile]{Timing Error Model}
 
A generic timing offset error model can be expressed as:
\begin{equation}
 r(t)=\sum\limits_{n}x(n)h(t-\tau(t)-nT_s)+v(t)
\end{equation}
where:
\begin{itemize}
 \item $x(n)$ is the transmitted symbol
 \item $h(t)$ is the pulse shape of the transmit filter
 \item $\tau(t)$ is the fractional delay offset
 \item $T_s$ is the sampling period
 \item $n$ is the sampling index
 \item $v(t)$ is the additive channel noise
\end{itemize}


\end{frame}


%%%%%%%%%%%%%%%%%%%%%%%%%%%%%%%%%%%%%%%%%%%%%%%%%%%%%%%%%%

\begin{frame}[fragile]{Timing Error Model -- Post-Reception}

Once the receive matched filter is applied, we get:
\begin{equation}
 y(t)=\sum\limits_{n}x(n)h_A(t-\tau(t)-nT_s)+v_h(t)
\end{equation}
where:
\begin{itemize}
 \item $h_A(t)=h(t)\ast{\bar{h}(-t)}$ $\rightarrow$ autocorrelation of transmit filter with its conjugate
 \item $v_h(t)$ is shaped noise
\end{itemize}

 
\end{frame}

%%%%%%%%%%%%%%%%%%%%%%%%%%%%%%%%%%%%%%%%%%%%%%%%%%%%%%%%%%

\begin{frame}[fragile]{Timing Error Model Illustration}

How does timing error look like between the transmitter and receiver when two free-running clocks are used?
 
\end{frame}


%%%%%%%%%%%%%%%%%%%%%%%%%%%%%%%%%%%%%%%%%%%%%%%%%%%%%%%%%%


\end{document}
